\chapter{Matchings and Flows}

\begin{center}\itshape
There is always a way to marry everybody so that no man and woman would both rather be married to each other.\\
\upshape(Gale and Shapley, 1962)
\end{center}

In the 1940s, American hospitals fought over medical graduates. Each wanted the best interns, so each made its offers a
little earlier than its rivals; by 1945 students were signing contracts two years before they graduated, before anyone could
tell how good they were. Deadlines were then imposed, and the market collapsed the other way: offers came with a few hours to
answer, students on a waiting list held out for their first choice while other hospitals' offers expired, and telephones rang
all night. In 1952 the hospitals gave up and created a clearing house. Students and hospitals sent in their ranked lists, and
an algorithm produced the assignment. Remarkably, participation was voluntary, and almost everybody accepted the result.

Ten years later David Gale and Lloyd Shapley, who had never heard of the clearing house, published a six-page paper on
``college admissions and the stability of marriage''. They asked for an assignment that nobody would want to break: no student
and hospital should both prefer each other to what they got. They proved that such a \emph{stable} assignment always exists, gave
a simple algorithm to find it, and showed that the algorithm is the best possible for the side that makes the offers. In 1984 Alvin
Roth discovered that the 1952 algorithm was exactly theirs: the doctors had found it by trial and error, and stability was why
the market held together. Roth and Shapley received the 2012 Nobel prize in economics for this theory; the same algorithm now
assigns pupils to schools in New York and Boston, and Roth's related designs organise kidney exchanges.

This chapter is about pairing things up. It starts with the purest form of the question, the matchings of a graph, and
Berge's principle that a matching can be improved exactly when an augmenting path exists. From it come Hall's marriage
theorem, K\H onig's theorem and a second proof of Dilworth's theorem from Chapter~2. Then come networks and flows, where
Menger's theorem of Chapter~5 becomes the max-flow min-cut theorem; Birkhoff's theorem on doubly stochastic matrices; Tutte's
theorem for graphs that are not bipartite; and finally the stable marriages that opened the chapter.

\begin{think}
Four students must share two double rooms. Each ranks the three others. Is there always a pairing in which no two students
in different rooms prefer each other to their roommates? And if there are two boys and two girls, and only mixed rooms are
allowed, with each student ranking the two others of the other sex?
\end{think}

\section{Matchings and augmenting paths}

A \emph{matching} of a graph $G$ is a set $M$ of edges no two of which share a vertex. The vertices covered by $M$ are \emph{matched},
the others are \emph{free}, and the matching is \emph{perfect} if it covers every vertex. A matching is \emph{maximum} if no matching has more
edges; this is stronger than \emph{maximal} (no edge can be added): in a path with three edges, the middle edge is a maximal
matching but not a maximum one. We write $\nu(G)$ for the size of a maximum matching.

An \emph{$M$-alternating path} is a path whose edges are alternately in $M$ and not in $M$; it is \emph{$M$-augmenting} if it joins two free
vertices (it then starts and ends with edges not in $M$, and has one more edge outside $M$ than inside). Swapping the two kinds
of edges along an augmenting path $P$ gives the matching $M\triangle P$, with one edge more (Figure~\ref{fig:augmenting}).

\begin{figure}[ht]
\centering
\fig{ch06-augmenting}
\caption{A matching of size $2$ (thick) and an augmenting path $x_1y_1x_2y_2x_3y_3$ (dashed), whose ends $x_1$ and $y_3$ are free.
Swapping matched and unmatched edges along the path gives a matching of size $3$.}
\label{fig:augmenting}
\end{figure}

\begin{theorem}[Berge, 1957]\label{thm:berge}
A matching $M$ of a graph $G$ is maximum if and only if there is no $M$-augmenting path.
\end{theorem}
\begin{proof}
If there is an augmenting path $P$, then $M\triangle P$ is a larger matching. Conversely, suppose $M'$ is a matching with $|M'|>|M|$, and
look at the edges of $M\triangle M'$. Each vertex meets at most one edge of $M$ and one of $M'$, so in the graph formed by $M\triangle M'$ every
vertex has degree at most $2$, and the components are paths and cycles whose edges alternate between $M$ and $M'$. An alternating
cycle has as many edges of $M$ as of $M'$. Since $M\triangle M'$ has more edges of $M'$ than of $M$, some component is a path with more
edges of $M'$ than of $M$: it starts and ends with edges of $M'$, and its two end vertices meet no edge of $M$ in $M\triangle M'$. They meet no edge
of $M\cap M'$ either, since they already meet an edge of $M'$. So they are free for $M$, and the path is $M$-augmenting.
\end{proof}

Berge's theorem is an algorithm in waiting: start with any matching, look for an augmenting path, swap, repeat. At most
$\frac n2$ rounds are needed. The difficulty is to \emph{find} an augmenting path. In a bipartite graph it is easy, by a search
that alternates between the two sides, as we shall see. In general graphs, odd cycles make the search loop back on
itself; Jack Edmonds solved this in 1965 by shrinking the odd cycles (``blossoms''), in a paper that also proposed
polynomial time as the definition of an efficient algorithm. Tome~C2 tells that story.

\section{Hall's marriage theorem and K\H onig's theorem}

From now until Section~\ref{sec:tutte}, $G$ is a bipartite graph with sides $X$ and $Y$: every edge joins $X$ to $Y$. For $S\subseteq X$, the
\emph{neighbourhood} $N(S)$ is the set of vertices of $Y$ adjacent to some vertex of $S$. If $M$ matches every vertex of $X$, the partners of $S$ are
$|S|$ distinct vertices of $N(S)$, so $|N(S)|\ge|S|$. This obvious necessary condition is enough.

\begin{theorem}[Hall's marriage theorem, 1935]\label{thm:hall}
A bipartite graph with sides $X$ and $Y$ has a matching covering $X$ if and only if
$|N(S)|\ge|S|$ for every $S\subseteq X$.
\end{theorem}

The traditional reading: $X$ is a set of women, $Y$ a set of men, edges join women to the men they would accept; everybody can be
married if and only if every group of $k$ women knows at least $k$ acceptable men. The theorem is so central that it is worth
proving three times, by three different philosophies.

\begin{proof}[First proof (Halmos and Vaughan, 1950)]
Induction on $|X|$; the case $|X|\le1$ is clear. Suppose first that $|N(S)|\ge|S|+1$ for every non-empty $S\subsetneq X$. Match any
vertex $x$ to any neighbour $y$, and delete both. In the remaining graph every non-empty $S\subseteq X\setminus\{x\}$ has lost at most the
neighbour $y$, so it still has at least $|S|$ neighbours, and induction matches $X\setminus\{x\}$.

Otherwise some non-empty $S_0\subsetneq X$ has $|N(S_0)|=|S_0|$. Hall's condition holds in the subgraph $G_0$ induced by $S_0\cup N(S_0)$
(the neighbours of a subset of $S_0$ are in $N(S_0)$), and $|S_0|<|X|$, so induction matches $S_0$ into $N(S_0)$. In the subgraph $G_1$ induced by
$(X\setminus S_0)\cup(Y\setminus N(S_0))$, a set $T\subseteq X\setminus S_0$ has neighbourhood $N(T\cup S_0)\setminus N(S_0)$, of size at least
$|T|+|S_0|-|S_0|=|T|$. So induction matches $X\setminus S_0$ in $G_1$, and the two matchings together cover $X$.
\end{proof}

\begin{proof}[Second proof (augmenting paths)]
Let $M$ be a maximum matching and suppose that a vertex $x\in X$ is free. Let $Z$ be the set of vertices reachable from $x$ by
$M$-alternating paths starting at $x$, and let $S=Z\cap X$ and $T=Z\cap Y$. Every $y\in T$ is matched, otherwise the alternating path from $x$
to $y$ would be augmenting, contradicting Berge's theorem; and the partner of $y$ is in $S$, since the path can be continued along
the matching edge. Conversely every vertex of $S$ other than $x$ is reached through its matching edge, so its partner is in $T$.
Hence $M$ pairs $T$ with $S\setminus\{x\}$, and $|S|=|T|+1$. Finally $N(S)=T$: a neighbour $y$ of $s\in S$ is reached by going to $s$ (along a path that
ends with a matching edge, or the empty path if $s=x$) and then along the edge $sy$, which is not in $M$ unless $y$ is the partner of $s$,
itself in $T$. So $|N(S)|=|S|-1$, and Hall's condition fails.
\end{proof}

The third proof, from the max-flow min-cut theorem, comes in Section~\ref{sec:flows}. The second proof is the efficient one: if
Hall's condition holds, a search from a free vertex finds an augmenting path; if it fails, the same search returns a set $S$ with
too few neighbours, a certificate that no matching exists. One computation, two possible answers, each checkable at a glance.

\begin{corollary}[Frobenius, 1917; K\H onig, 1916]\label{cor:regular}
Let $k\ge1$. A $k$-regular bipartite multigraph has a perfect matching, and its edges can be partitioned into $k$ perfect matchings.
\end{corollary}
\begin{proof}
Counting edges, $k|X|=k|Y|$, so $|X|=|Y|$. For $S\subseteq X$, the $k|S|$ edges leaving $S$ all end in $N(S)$, which receives $k|N(S)|$ edges in total;
so $k|S|\le k|N(S)|$ and Hall's condition holds (parallel edges do not matter for it). A perfect matching exists; removing it leaves a
$(k-1)$-regular bipartite multigraph, and induction on $k$ gives the partition.
\end{proof}

A \emph{vertex cover} is a set $K$ of vertices meeting every edge. Each edge of a matching needs its own vertex in $K$, so
$\nu(G)\le|K|$ for every matching and every cover. In the triangle, $\nu=1$ and the smallest cover has $2$ vertices. In a bipartite
graph there is never a gap.

\begin{theorem}[K\H onig, 1931; Egerv\'ary, 1931]\label{thm:konig}
In a bipartite graph, the maximum size of a matching equals the minimum size of a vertex cover.
\end{theorem}
\begin{proof}
Let $M$ be a maximum matching, let $U$ be the set of free vertices of $X$, and let $Z$ be the set of vertices reachable from $U$ by
$M$-alternating paths starting in $U$ (including $U$ itself). Put
\[
K=(X\setminus Z)\cup(Y\cap Z).
\]
$K$ is a cover. Let $xy$ be an edge, $x\in X$. If $x\notin Z$, then $x\in K$. If $x\in Z$ and $xy\notin M$, the alternating path to $x$ ends with an edge of $M$ (or
is empty, if $x\in U$), so it continues along $xy$, and $y\in Z$. If $x\in Z$ and $xy\in M$, then $x\notin U$, the path reaching $x$ arrives through its
matching edge, and so $y\in Z$ as well. In each case $xy$ meets $K$.

$|K|=|M|$. Every $y\in Y\cap Z$ is matched (otherwise there is an augmenting path), and its partner is in $Z$. Every $x\in X\setminus Z$ is matched, since
$U\subseteq Z$, and its partner is not in $Z$ (otherwise $x$ would be reached through it). So every vertex of $K$ is matched, and each edge $xy$ of
$M$ has exactly one end in $K$: $y$ if $y\in Z$ (then $x\in Z$ too), $x$ if $y\notin Z$ (then $x\notin Z$, since a vertex $x\in Z$ matched to $y$
is reached through $y$). Hence $|K|=|M|$, and since every cover has at least $\nu(G)$ vertices, $K$ is a minimum cover.
\end{proof}

Again the proof is a certificate: the matching $M$ and the cover $K$, of the same size, prove each other optimal. In the
language of Tome~C5 this is the strong duality of a pair of linear programs, and the theorem is the first example of a
\emph{min-max} theorem, of which this chapter contains half a dozen. Here is one that we met before, with a proof due to Delbert Fulkerson (1956).

\begin{proof}[Second proof of Theorem~\ref{thm:dilworth}]
Let $P$ be a finite poset with $n$ elements. Build a bipartite graph with two copies $X=\{x^-:x\in P\}$ and $Y=\{y^+:y\in P\}$ of $P$, and an edge
$x^-y^+$ whenever $x<y$. A matching $M$ is a choice, for some elements $x$, of a ``successor'' $y>x$, with no element chosen twice as a
successor. Following successors, the elements fall into chains (each element has at most one successor and at most one predecessor),
and each edge of $M$ glues two of the $n$ singletons together: the number of chains is $n-|M|$.

Now take a maximum matching $M$ and, by K\H onig's theorem, a cover $K$ with $|K|=|M|$. Let $A$ be the set of $x\in P$ such that neither
$x^-$ nor $x^+$ is in $K$; then $|A|\ge n-|K|$. If $x<y$ were both in $A$, the edge $x^-y^+$ would not be covered: so $A$ is an antichain. Hence
$P$ has a partition into $n-|M|\le|A|$ chains, and an antichain of size $|A|$; since an antichain meets each chain at most once, the
minimum number of chains equals the maximum size of an antichain.
\end{proof}

This proof is an algorithm: a maximum bipartite matching takes polynomial time (Tome~C2), and the chains and the antichain
come out of it. The same construction on a directed acyclic graph computes the smallest number of paths covering its
vertices (Exercise~\ref{ex:pathcover}), a problem met in scheduling.

\section{Flows in networks}\label{sec:flows}

In 1955 two analysts of the RAND Corporation, T.~E.~Harris and F.~S.~Ross, wrote a secret report for the US Air Force on the
rail network linking the Soviet Union to Eastern Europe: $44$ regions, $105$ links, each with a capacity in tons per day. How much
could the network carry, and which links should a bomber cut to stop it most cheaply? Their colleagues Lester Ford and Delbert
Fulkerson turned the question into a theorem the next year. The report was declassified in 1999.

A \emph{network} is a directed graph $D=(V,A)$ with two vertices, a \emph{source} $s$ and a \emph{sink} $t\ne s$, and a \emph{capacity} $c(a)\ge0$
on each arc. A \emph{flow} is a function $f:A\to\R$ with
\begin{itemize}
\item $0\le f(a)\le c(a)$ for every arc (capacity), and
\item $f^{\mathrm{in}}(v)=f^{\mathrm{out}}(v)$ for every $v\ne s,t$ (conservation),
\end{itemize}
where $f^{\mathrm{out}}(v)$ and $f^{\mathrm{in}}(v)$ are the sums of $f$ over the arcs leaving and entering $v$. Its \emph{value} is
$|f|=f^{\mathrm{out}}(s)-f^{\mathrm{in}}(s)$. A \emph{cut} is a set $S$ of vertices with $s\in S$, $t\notin S$; its \emph{capacity} $c(S)$ is the sum
of the capacities of the arcs leaving $S$.

\begin{lemma}\label{lem:weakduality}
For every flow $f$ and every cut $S$, $|f|=f^{\mathrm{out}}(S)-f^{\mathrm{in}}(S)\le c(S)$, where $f^{\mathrm{out}}(S)$ and $f^{\mathrm{in}}(S)$ are the
sums of $f$ over the arcs leaving and entering $S$.
\end{lemma}
\begin{proof}
Sum $f^{\mathrm{out}}(v)-f^{\mathrm{in}}(v)$ over $v\in S$. Only $s$ contributes, giving $|f|$. On the other hand an arc with both ends in $S$ contributes
$+f(a)$ at its tail and $-f(a)$ at its head; what remains is $f^{\mathrm{out}}(S)-f^{\mathrm{in}}(S)$, which is at most
$f^{\mathrm{out}}(S)\le c(S)$.
\end{proof}

Given a flow $f$, an \emph{augmenting path} is a path from $s$ to $t$ in the underlying undirected graph, in which every arc traversed
forwards has $f(a)<c(a)$ and every arc traversed backwards has $f(a)>0$. If $\delta>0$ is the smallest of the slacks $c(a)-f(a)$ (forward
arcs) and $f(a)$ (backward arcs), adding $\delta$ on forward arcs and subtracting $\delta$ on backward arcs gives a flow: capacities hold by the
choice of $\delta$, and at each inner vertex of the path the two changes cancel (check the four cases: in and out, both forward,
both backward, or one of each). The value grows by $\delta$, since the first arc leaves $s$ forwards or enters it backwards.

\begin{theorem}[Max-flow min-cut; Ford and Fulkerson, 1956; Elias, Feinstein and Shannon, 1956]\label{thm:maxflow}
In every network:
\begin{enumerate}
\item a flow is maximum if and only if it has no augmenting path, and then the set $S$ of vertices reachable from $s$ by
augmenting-type paths is a cut with $c(S)=|f|$;
\item the maximum value of a flow equals the minimum capacity of a cut;
\item if all capacities are integers, there is a maximum flow with integer values, and the augmenting path algorithm started
from the zero flow finds one.
\end{enumerate}
\end{theorem}
\begin{proof}
(1) An augmenting path increases the value, so a maximum flow has none. Conversely, let $f$ have no augmenting path, and let $S$
be the set of vertices $v$ that can be reached from $s$ by a path of the same type as an augmenting path (forward arcs not full,
backward arcs not empty). Then $s\in S$ and $t\notin S$. An arc $uv$ leaving $S$ is full, $f(uv)=c(uv)$, otherwise $v$ would be reachable;
an arc $vu$ entering $S$ is empty, otherwise $v$ would be reachable backwards. By Lemma~\ref{lem:weakduality}, $|f|=c(S)$; and since every
flow has value at most $c(S)$, $f$ is maximum and $S$ is a minimum cut.

(2) The flows form a closed and bounded subset of $\R^A$ (defined by finitely many linear equations and inequalities), which is
non-empty (the zero flow) and compact. The value is a continuous function, so a maximum flow exists; by (1) its value is the
capacity of some cut, which is minimum by Lemma~\ref{lem:weakduality}.

(3) Start from the zero flow. If capacities and flow are integers, the slack $\delta$ of an augmenting path is a positive integer,
so the new flow is integral and its value grows by at least $1$. The value is at most $c(\{s\})$, so after at most $c(\{s\})$
augmentations no augmenting path remains, and by (1) the integral flow obtained is maximum.
\end{proof}

\begin{figure}[ht]
\centering
\fig{ch06-flow}
\caption{A maximum flow, each arc labelled flow/capacity, of value $19$. The cut $S=\{s,c\}$ (dashed) is crossed forwards by
$s\to a$ and $c\to d$, both full, of total capacity $10+9=19$, and backwards by $a\to c$, which is empty: the cut certifies that no
flow does better.}
\label{fig:flow}
\end{figure}

The compactness argument in (2) is there for a reason: with irrational capacities and an unlucky choice of augmenting paths,
the algorithm can run forever, and its values may even converge to less than the maximum (Ford and Fulkerson gave such an
example; Uri Zwick found the smallest ones in 1995). Choosing a \emph{shortest} augmenting path each time, as Edmonds and Karp proposed
in 1972, makes the algorithm terminate after at most $|V||A|$ augmentations whatever the capacities; Tome~C2 proves this and
goes on to the modern algorithms.

Part (3), the \emph{integrality theorem}, is what makes flows a combinatorial tool: whenever a problem asks for an integer
assignment subject to capacity constraints, one may solve the continuous problem and get the integer answer for free. To read
paths in a flow we need one more lemma.

\begin{lemma}\label{lem:decomp}
If all capacities are $1$ and $f$ is an integral flow of value $k\ge0$, the arcs with $f(a)=1$ contain $k$ arc-disjoint directed paths from $s$ to $t$.
\end{lemma}
\begin{proof}
Induction on $k$, the case $k=0$ being empty. Let $H$ be the set of arcs with $f(a)=1$. In $H$, every vertex other than $s$ and $t$ has as
many arcs in as out, and $s$ has $k\ge1$ more arcs out than in. Walk from $s$ along arcs of $H$, never using an arc twice. Each time the walk enters a
vertex $v\ne s,t$, it has used one more arc into $v$ than out of it, so an unused arc leaves $v$; each time it is at $s$, it has used as many arcs out of $s$ as
into it, so again an unused arc leaves. As $H$ is finite, the walk stops, hence at $t$. By Lemma~\ref{lem:walkpath}, which works
the same way for directed walks, it contains a directed path $Q$ from $s$ to $t$. Setting $f=0$ on the arcs of $Q$ gives an integral flow of
value $k-1$ (conservation is kept at the inner vertices of $Q$), whose arcs contain $k-1$ arc-disjoint paths by induction; with $Q$, that makes $k$.
\end{proof}

\begin{corollary}[Menger's theorem for arcs]\label{cor:edgemenger}
In a directed graph, the maximum number of arc-disjoint directed paths from $s$ to $t$ equals the minimum number of arcs whose
removal leaves no directed path from $s$ to $t$.
\end{corollary}
\begin{proof}
Give every arc capacity $1$. A set $F$ of arcs meeting every $s$--$t$ path contains all the arcs leaving the cut $S$ formed by the
vertices still reachable from $s$ in $D-F$, so $|F|\ge c(S)$; conversely the arcs leaving any cut meet every $s$--$t$ path. So the
minimum in the statement is the minimum capacity of a cut. By Theorem~\ref{thm:maxflow} there is an integral flow of that value, and
Lemma~\ref{lem:decomp} turns it into as many arc-disjoint paths. Fewer cannot be: each path uses an arc of each such $F$.
\end{proof}

For an undirected graph, replace every edge by two opposite arcs. If two of the paths obtained use the arcs $uv$ and $vu$ of the
same edge, exchange their tails after $u$ and $v$ and drop the two arcs: the paths become walks in the same graph, from which
paths can be extracted, and the number of edges used goes down. So the edge version of Menger's theorem holds for graphs too, and
splitting each vertex into an arc of capacity $1$ gives the vertex version, Theorem~\ref{thm:menger}, once more. Hall's theorem is
also a cut in disguise.

\begin{proof}[Third proof of Hall's theorem]
Add a source $s$ with an arc $s\to x$ of capacity $1$ to each $x\in X$, a sink $t$ with an arc $y\to t$ of capacity $1$ from each $y\in Y$, and
orient each edge from $X$ to $Y$ with capacity $|X|+1$. An integral flow of value $|X|$ sends one unit through each $x\in X$, along one edge,
to distinct vertices of $Y$: a matching covering $X$. By Theorem~\ref{thm:maxflow} it suffices that every cut has capacity at least
$|X|$. Let $S$ be a cut and $A=S\cap X$, $B=S\cap Y$. If an edge goes from $A$ to $Y\setminus B$, then $c(S)\ge|X|+1$. Otherwise $N(A)\subseteq B$,
and the arcs leaving $S$ include the arcs $s\to x$ for $x\in X\setminus A$ and $y\to t$ for $y\in B$:
$c(S)\ge|X\setminus A|+|B|\ge|X|-|A|+|N(A)|\ge|X|$.
\end{proof}

Flows model far more than pipes and railways. Exercise~\ref{ex:baseball} decides whether a baseball team can still finish first,
and Exercise~\ref{ex:closure} chooses the most profitable set of projects subject to prerequisites; the same minimum cuts separate
an object from its background in image segmentation.

\section{Doubly stochastic matrices}

A square matrix with non-negative entries is \emph{doubly stochastic} if every row and every column sums to $1$; it is the transition
matrix of a Markov chain (Tome~3) whose uniform distribution is stationary. The permutation matrices are doubly stochastic, and so
is any average of them, for instance $\frac1n J$, the matrix with all entries $\frac1n$. The converse is the content of the following theorem.

\begin{theorem}[Birkhoff, 1946; von Neumann, 1953]\label{thm:birkhoff}
Every doubly stochastic matrix is a convex combination of permutation matrices. If it has $p$ positive entries, at most
$p-n+1$ permutation matrices are needed.
\end{theorem}
\begin{proof}
Let $A$ be doubly stochastic of size $n$, and $G$ the bipartite graph with rows and columns as sides and an edge $ij$ whenever $a_{ij}>0$.
For a set $S$ of rows, the entries of these rows add up to $|S|$, and they all lie in the columns of $N(S)$, whose entries add up to $|N(S)|$;
so $|S|\le|N(S)|$. By Hall's theorem $G$ has a perfect matching, that is a permutation $\sigma$ with $a_{i\sigma(i)}>0$ for all $i$. Let $P$ be its
permutation matrix and $\lambda=\min_ia_{i\sigma(i)}$, so $0<\lambda\le1$.

We induct on $p\ge n$ (each row has a positive entry). If $p=n$, each row has a single positive entry, equal to $1$, so $A=P$ and
$1=p-n+1$ matrix suffices. If $p>n$, then $A\ne P$, so $\lambda<1$, and
\[
A'=\frac{A-\lambda P}{1-\lambda}
\]
has non-negative entries and row and column sums $\frac{1-\lambda}{1-\lambda}=1$: it is doubly stochastic, with at least one positive entry fewer
(where the minimum $\lambda$ is attained). By induction $A'$ is a convex combination of at most $(p-1)-n+1$ permutation matrices,
and $A=\lambda P+(1-\lambda)A'$ of at most $p-n+1$.
\end{proof}

Since $p\le n^2$, at most $n^2-n+1$ permutation matrices are needed; a dimension count (Carath\'eodory's theorem, Tome~C5) brings this down
to $(n-1)^2+1$ (Marcus and Ree, 1959).

\begin{beauty}
The assignment problem asks for a permutation $\sigma$ minimising a cost $\sum_ic_{i\sigma(i)}$: $n$ workers, $n$ jobs, $n!$ possibilities.
Relax it to doubly stochastic matrices, minimising the linear function $\sum_{i,j}c_{ij}a_{ij}$ over a convex set: a linear program,
solvable fast. Nothing is lost. By Birkhoff's theorem the cost of $A=\sum_k\lambda_kP_k$ is the average $\sum_k\lambda_k\,\mathrm{cost}(P_k)$,
so some permutation is at least as cheap as $A$. A problem that looks like a search through $n!$ cases is secretly convex.
\end{beauty}

\section{Perfect matchings in general graphs}\label{sec:tutte}

Without bipartiteness Hall's condition is meaningless, and odd cycles get in the way: a triangle has no perfect matching
although every vertex has two neighbours. The obstruction is an odd component. If $G-S$ (remove the vertices of $S$) has a
component with an odd number of vertices, a perfect matching must match one of its vertices with a vertex of $S$. Write $o(H)$ for
the number of components of $H$ with an odd number of vertices. If $o(G-S)>|S|$, there are not enough vertices in $S$ to go round
(Figure~\ref{fig:sylvester}).

\begin{figure}[ht]
\centering
\fig{ch06-sylvester}
\caption{A graph in which every vertex has degree $3$, but with no perfect matching: removing the vertex $s$ leaves three
components of $5$ vertices each, and $s$ can be matched into only one of them. Its three edges at $s$ are bridges.}
\label{fig:sylvester}
\end{figure}

\begin{theorem}[Tutte, 1947]\label{thm:tutte}
A graph $G$ has a perfect matching if and only if $o(G-S)\le|S|$ for every set $S$ of vertices.
\end{theorem}
\begin{proof}[Proof (Lov\'asz, 1975)]
Necessity was explained above: each odd component of $G-S$ sends a matching edge to its own vertex of $S$. For sufficiency,
suppose that $G$ satisfies Tutte's condition but has no perfect matching. With $S=\varnothing$, the condition says that $G$ has no odd component,
so $n=|V|$ is even. Add edges to $G$, one at a time, as long as there is still no perfect matching; we reach a graph $G^*$ on the same vertices,
without perfect matching, such that adding any missing edge creates one. Adding an edge never increases $o(G-S)$ (it may merge two
components: two odd ones become one even one, an odd and an even one an odd one), so $G^*$ still satisfies Tutte's condition. We now
work in $G^*$, and call $S$ the set of vertices adjacent to all other vertices.

\emph{If every component of $G^*-S$ is a complete graph}, we build a perfect matching, a contradiction. Match the vertices of each
even component among themselves. In each odd component, match all vertices but one among themselves, and match the remaining
vertex with a vertex of $S$, a different one for each component; this is possible since $o(G^*-S)\le|S|$ and the vertices of $S$ are
adjacent to everything. The vertices of $S$ left over are pairwise adjacent, and their number $|S|-o(G^*-S)$ is even, because
$n=|S|+\sum_C|C|\equiv|S|+o(G^*-S)\pmod2$ and $n$ is even. Match them in pairs.

\emph{Otherwise} some component of $G^*-S$ is not complete. A connected graph that is not complete contains vertices $a,b,a'$ with $ab$ and
$ba'$ edges but $aa'$ not an edge (take two non-adjacent vertices at distance $2$). Since $b\notin S$, some vertex $c$ is not adjacent to $b$. By the
maximality of $G^*$, the graph $G^*+aa'$ has a perfect matching $M_1$, which contains $aa'$, and $G^*+bc$ has a perfect matching $M_2$, which contains $bc$.
As in Berge's proof, $M_1\triangle M_2$ is a disjoint union of cycles alternating between $M_1$ and $M_2$ (both matchings are perfect, so
there are no paths). Let $C$ be the cycle containing $bc$. Note that if $a$ or $a'$ is on $C$, then so is the edge $aa'$, the $M_1$-edge at $a$.
\begin{itemize}
\item If $aa'$ is not on $C$, take $M_1$ on the vertices of $C$ and $M_2$ elsewhere. This is a perfect matching which avoids $bc$ (an $M_2$-edge
of $C$) and $aa'$ (not on $C$, and not in $M_2$), so it lies in $G^*$: a contradiction.
\item If $aa'$ is on $C$, walk along $C$ from $b$, starting with the $M_1$-edge at $b$, until the first vertex $d\in\{a,a'\}$, and let $d'$ be the
other one. This path $P$ starts with an $M_1$-edge and ends with an $M_2$-edge (the $M_1$-edge at $d$ is $dd'$, and $d'$ comes later), and it
does not contain $c$, the last vertex before the walk returns to $b$ (and $c\ne a,a'$, as $c$ is not adjacent to $b$). Take the $M_2$-edges of
$P$, which cover every vertex of $P$ except $b$; the edge $bd'$, which is in $G^*$ since $b$ is adjacent to $a$ and $a'$; and the $M_1$-edges with
no end on $P\cup\{d'\}$, which cover all other vertices (the $M_1$-partners of the vertices of $P$ other than $d$ are on $P$, and that of $d$ is
$d'$). This is a perfect matching of $G^*$, using neither $bc$ nor $aa'$: a contradiction again.
\end{itemize}
\end{proof}

Tutte's original proof used Pfaffians, the square roots of determinants of skew-symmetric matrices, and the idea survives in a
fast randomised test for perfect matchings by Lov\'asz (1979) that Tome~C3 presents. Here is the oldest theorem of the subject,
of which Figure~\ref{fig:sylvester} shows that its hypothesis cannot be dropped.

\begin{corollary}[Petersen, 1891]\label{cor:petersen}
Every $3$-regular graph without a bridge has a perfect matching.
\end{corollary}
\begin{proof}
Let $S$ be a set of vertices and $C$ an odd component of $G-S$. The degrees in $C$ add up to $3|C|$, which is odd, and count twice the
edges inside $C$ plus the edges from $C$ to $S$ (there are no edges to other components). So the number $e(C,S)$ of edges from $C$ to $S$ is odd.
It is not $1$, since that edge would be a bridge (it would be the only edge between $C$ and the rest of $G$), so $e(C,S)\ge3$. The edges from
odd components to $S$ number at most $3|S|$, since each vertex of $S$ has degree $3$; hence $3\,o(G-S)\le3|S|$, which is Tutte's condition.
\end{proof}

\begin{remark}
The Gallai--Edmonds structure theorem (1963--65) goes further and describes all maximum matchings of a graph through a
canonical set $S$ attaining the maximum of $o(G-S)-|S|$; the Tutte--Berge formula $\nu(G)=\frac12\bigl(n-\max_S(o(G-S)-|S|)\bigr)$
follows from Theorem~\ref{thm:tutte} as Exercise~\ref{ex:konigore} follows from Hall's theorem. See Lov\'asz and Plummer's book.
\end{remark}

\section{Stable marriages}

Back to the doctors. There are $n$ \emph{proposers} and $n$ \emph{receivers}; each proposer ranks all receivers in a strict order, and
each receiver ranks all proposers. A perfect matching $M$ between the two sides is \emph{unstable} if some proposer $p$ and receiver $r$,
not matched together, both prefer each other to their partners in $M$; the pair $(p,r)$ is then \emph{blocking}. Otherwise $M$ is
\emph{stable}. Gale and Shapley's algorithm, called \emph{deferred acceptance}, runs as follows.
\begin{itemize}
\item While some proposer $p$ is free, $p$ proposes to the receiver $r$ it ranks highest among those it has not yet proposed to.
\item The receiver $r$ holds the proposal it prefers among the one it holds (if any) and the new one, and rejects the other; the
rejected proposer becomes free.
\end{itemize}
The word ``deferred'' is the key: a receiver never accepts definitively, so it can trade up until the end.

\begin{figure}[ht]
\centering
\fig{ch06-galeshapley}
\caption{Preferences of four proposers and four receivers, best first. Deferred acceptance gives every proposer the receiver in
bold in its list, and its outcome $Ax$, $Bz$, $Cw$, $Dy$ is stable. Matching every receiver with its first choice (bold on the right),
that is $Bw$, $Cx$, $Ay$, $Dz$, is stable too.}
\label{fig:galeshapley}
\end{figure}

On the preferences of Figure~\ref{fig:galeshapley}, $A$ proposes to $x$, $B$ to $y$ and $C$ to $z$, and all are held. Then $D$ proposes to $y$,
who prefers $D$ to $B$ and rejects $B$. Now $B$ tries $x$, who keeps $A$; then $z$, who prefers $B$ to $C$. Rejected, $C$ tries $w$, who is free.
After seven proposals everybody is matched: $Ax$, $Bz$, $Cw$, $Dy$.

\begin{theorem}[Gale and Shapley, 1962]\label{thm:galeshapley}
Deferred acceptance stops after at most $n^2$ proposals, and the matching it produces is stable. In particular a stable matching
always exists.
\end{theorem}
\begin{proof}
A proposer never proposes twice to the same receiver, so there are at most $n^2$ proposals. A receiver who has received a proposal
holds one from then on, and the proposals it holds get better and better (for it). A proposer $p$ is never rejected by every receiver: if
$p$ had been rejected by all $n$ receivers, all of them would have received a proposal and would hold one, from the $n-1$ other proposers, which
is impossible. So when the algorithm stops, no proposer is free, and the held proposals form a perfect matching $M$.

Suppose that $p$ prefers $r$ to its partner in $M$. Then $p$ proposed to $r$ before proposing to its final partner, and was rejected by $r$,
either at once or later, in favour of a proposer that $r$ preferred. Since the proposals held by $r$ only improve, $r$ prefers its final partner to
$p$. So $(p,r)$ is not blocking, and $M$ is stable.
\end{proof}

\begin{theorem}[Gale and Shapley, 1962; McVitie and Wilson, 1971]\label{thm:proposeroptimal}
Call a receiver \emph{achievable} for a proposer if some stable matching pairs them. The matching produced by deferred acceptance
gives every proposer its best achievable receiver, and every receiver its worst achievable proposer. In particular it does
not depend on the order in which free proposers are chosen.
\end{theorem}
\begin{proof}
We show that no proposer is ever rejected by an achievable receiver. Suppose the contrary, and consider the first time it
happens during the run: $r$ rejects $p$ (at once or later) in favour of $p'$, which $r$ prefers to $p$, although some stable matching $M$ pairs $p$
with $r$. In $M$, $p'$ has a partner $r'\ne r$. When $p'$ proposed to $r$, it had already been rejected by every receiver it ranks above $r$, and
all those rejections came before the first rejection by an achievable receiver; so none of these receivers is achievable for $p'$,
and $r'$, which is achievable for $p'$, is not among them: $p'$ prefers $r$ to $r'$. But $r$ prefers $p'$ to $p$, its partner in $M$. So $(p',r)$ blocks $M$,
a contradiction.

Hence every proposer, proposing in decreasing order of preference and never rejected by an achievable receiver, ends with a
receiver at least as good as all its achievable ones; and its final receiver is achievable, since the output is stable. So it gets
its best achievable receiver. Now let $M_0$ be the output, and suppose a receiver $r$ is matched with $p$ in $M_0$ and prefers $p$ to its partner
$p''$ in some stable matching $M$. In $M$, $p$ has a partner $r''\ne r$, achievable for $p$, so $p$ prefers $r$, its best achievable receiver, to $r''$.
Then $(p,r)$ blocks $M$, a contradiction: $r$ gets its worst achievable proposer. Since the output has been described without
reference to the order of proposals, it does not depend on that order.
\end{proof}

So the side that proposes wins, and wins as much as stability permits; in Figure~\ref{fig:galeshapley}, if the receivers propose,
every one of them gets its first choice, and three of the four proposers their last. When the residency match was redesigned in
1998, the students were made the proposing side. Proposers also have no reason to lie: misreporting one's preferences never helps a
proposer (Dubins and Freedman, 1981), whereas a receiver can gain by lying (Exercise~\ref{ex:lying}).

How good is the outcome for random preferences? Knuth showed that the expected number of proposals is about $n\ln n$, a coupon
collector in disguise; Pittel showed that the proposers get on average their receiver of rank about $\ln n$, and the receivers their
proposer of rank about $n/\ln n$. The notebook measures both.

Real matching markets have more structure. In the hospitals-residents problem, each hospital has a number of positions; cloning
each hospital into as many copies as it has positions, each with the same preferences, reduces it to the problem above, and
students may also declare some hospitals unacceptable. With couples who want to be placed together, stable matchings may fail to
exist, and deciding whether one exists is NP-complete (Ronn, 1990), a reason why the real match uses heuristics for them. And when
everybody is on the same side, as in the room-sharing question at the start of the chapter, a stable pairing may not exist at all
(Exercise~\ref{ex:roommates}); Irving gave in 1985 an algorithm that finds one when it exists.

\begin{beauty}
Stable matchings have the structure of a lattice. Given two stable matchings, give every proposer the better of its two partners:
the result is again a stable matching, in which every receiver gets the worse of its two (Exercise~\ref{ex:stablelattice}). The
proposer-optimal and receiver-optimal matchings are the top and bottom of this lattice, and every finite distributive lattice
arises this way (Blair, 1984). The same order-theoretic structure as in Chapter~2 governs a problem from economics.
\end{beauty}

\section{Exercises}

\begin{exercise}[$\star\star$]\label{ex:latin}
A \emph{$k\times n$ Latin rectangle} is a $k\times n$ array of symbols $1,\dots,n$ in which every row contains each symbol once and no column contains a
symbol twice. Show that if $k<n$, it can be extended to a $(k+1)\times n$ Latin rectangle, and so, row by row, to a Latin square.
\end{exercise}

\begin{exercise}[$\star$]\label{ex:cards}
A deck of $52$ cards is dealt into $13$ piles of $4$ cards. Show that one can pick one card from each pile so that the $13$ cards picked
contain one card of each rank, Ace to King.
\end{exercise}

\begin{exercise}[$\star\star$]\label{ex:konigore}
\emph{The deficiency version.} In a bipartite graph with sides $X$ and $Y$, show that
$\nu(G)=|X|-\max_{S\subseteq X}\bigl(|S|-|N(S)|\bigr)$ (the maximum is at least $0$, from $S=\varnothing$).
\end{exercise}

\begin{exercise}[$\star\star$]\label{ex:edgecolour}
\emph{K\H onig's edge colouring theorem, 1916.} Show that the edges of a bipartite graph of maximum degree $\Delta$ can be coloured with $\Delta$ colours so
that edges sharing a vertex get different colours. Deduce how to schedule a round of exams in the fewest time slots, when each exam involves one student
and one examiner.
\end{exercise}

\begin{exercise}[$\star$]\label{ex:pathcover}
Show that the vertices of a directed acyclic graph with $n$ vertices can be covered by $n-\nu$ vertex-disjoint directed paths, and not fewer,
where $\nu$ is the maximum size of a matching in the bipartite graph with an edge $u^-v^+$ for each arc $u\to v$.
\end{exercise}

\begin{exercise}[$\star\star$]\label{ex:baseball}
\emph{Baseball elimination (Schwartz, 1966).} Teams $1,\dots,m$ have $w_i$ wins so far, and $g_{ij}$ games left between $i$ and $j$; there are no draws. Team $k$ has $r_k$
games left; let $W=w_k+r_k$. Assume $W\ge w_i$ for all $i$. Build a network with a source $s$, a vertex for each pair $\{i,j\}$ of teams other than $k$, joined from $s$ by an arc
of capacity $g_{ij}$, arcs of capacity $g_{ij}$ from $\{i,j\}$ to $i$ and to $j$, and an arc of capacity $W-w_i$ from each team $i\ne k$ to the sink $t$. Show that team $k$ can still
end with at least as many wins as every other team if and only if the maximum flow saturates every arc leaving $s$.
\end{exercise}

\begin{exercise}[$\star\star\star$]\label{ex:closure}
\emph{Project selection.} Projects $1,\dots,m$ have profits $p_i\in\R$ (costs if negative), and some projects require others: choosing $i$
forces choosing $j$ when $i\to j$ is a requirement. Build a network with an arc $s\to i$ of capacity $p_i$ for $p_i>0$, an arc $i\to t$ of capacity
$-p_i$ for $p_i<0$, and an arc $i\to j$ of capacity $+\infty$ for every requirement (or $1+\sum_i|p_i|$). Show that the largest total profit of an
allowed set of projects is $\sum_{p_i>0}p_i$ minus the minimum capacity of a cut.
\end{exercise}

\begin{exercise}[$\star\star$]\label{ex:birkhoffext}
Show that the permutation matrices are exactly the doubly stochastic matrices that are not the midpoint $\frac12(B+C)$ of two different doubly
stochastic matrices (the extreme points of the Birkhoff polytope).
\end{exercise}

\begin{exercise}[$\star\star$]\label{ex:errera}
\emph{Errera, 1922.} Show that the graph of Figure~\ref{fig:sylvester} is $3$-regular without a perfect matching, and that a $3$-regular graph with at most
two bridges has a perfect matching.
\end{exercise}

\begin{exercise}[$\star$]\label{ex:roommates}
Four students $a,b,c,d$ share two double rooms. Student $a$ prefers $b$ to $c$ to $d$, student $b$ prefers $c$ to $a$ to $d$, student $c$ prefers $a$ to $b$ to $d$,
and $d$ has any preferences. Show that every pairing has a blocking pair. Then answer the mixed-rooms question at the start of the chapter.
\end{exercise}

\begin{exercise}[$\star\star$]\label{ex:stablelattice}
Let $M$ and $M'$ be stable matchings. Give each proposer the better of its partners in $M$ and $M'$. Show that the result is a perfect matching,
that it is stable, and that each receiver gets the worse of its two partners.
\end{exercise}

\begin{exercise}[$\star\star$]\label{ex:lying}
Proposers $A,B,C$ rank the receivers $x\,y\,z$, $x\,z\,y$ and $y\,x\,z$ respectively (best first); receivers $x,y,z$ rank the proposers $C\,A\,B$, $A\,B\,C$ and $A\,B\,C$.
Run deferred acceptance. Then show that $x$, by reporting $C\,B\,A$ instead of its true ranking, ends with a partner it truly prefers.
\end{exercise}

\section{Hints}

\begin{hint}{ex:latin}
Join each column to the symbols missing from it, and count degrees.
\end{hint}
\begin{hint}{ex:cards}
A bipartite multigraph between piles and ranks, one edge per card.
\end{hint}
\begin{hint}{ex:konigore}
Add $d$ new vertices to $Y$, adjacent to all of $X$.
\end{hint}
\begin{hint}{ex:edgecolour}
Embed the graph in a $\Delta$-regular bipartite multigraph.
\end{hint}
\begin{hint}{ex:pathcover}
Read a matching as a choice of successors, as in the second proof of Dilworth's theorem.
\end{hint}
\begin{hint}{ex:baseball}
A unit of flow through $\{i,j\}$ and $i$ is a game between $i$ and $j$ won by $i$. Use integrality.
\end{hint}
\begin{hint}{ex:closure}
A cut of finite capacity is $\{s\}$ plus an allowed set of projects.
\end{hint}
\begin{hint}{ex:birkhoffext}
For one direction use Theorem~\ref{thm:birkhoff}; for the other, look at the entries equal to $0$ and $1$.
\end{hint}
\begin{hint}{ex:errera}
As in Corollary~\ref{cor:petersen}; parity of $n$ helps at the end.
\end{hint}
\begin{hint}{ex:roommates}
By symmetry, it is enough to examine who shares a room with $d$.
\end{hint}
\begin{hint}{ex:stablelattice}
If two proposers got the same receiver, one got it from $M$ and the other from $M'$; use the stability of both.
\end{hint}
\begin{hint}{ex:lying}
Trace the proposals one by one in both runs.
\end{hint}

\section{Solutions}

\begin{solution}{ex:latin}
\idea Adding a row means choosing, for each column, a missing symbol, all different: a perfect matching between columns and symbols.
\writeup Join column $j$ to each symbol not yet in column $j$. Each column misses $n-k$ symbols. Each symbol appears once in each of the $k$ rows, in $k$ different
columns, so it is missing from $n-k$ columns. The bipartite graph is $(n-k)$-regular with $n-k\ge1$, so by Corollary~\ref{cor:regular} it has a perfect
matching, which is a new row: each symbol once, and no repetition in a column.
\end{solution}

\begin{solution}{ex:cards}
\idea Each pile has $4$ cards, each rank has $4$ cards: this is regularity.
\writeup Join each pile to the rank of each of its cards, with multiplicity. Every pile has degree $4$, and every rank has degree $4$, since its four
cards are somewhere. By Corollary~\ref{cor:regular} there is a perfect matching: one card per pile, one per rank.
\end{solution}

\begin{solution}{ex:konigore}
\idea Each deficient set leaves vertices unmatched; conversely, a few universal extra partners repair Hall's condition.
\writeup Let $d=\max_S(|S|-|N(S)|)\ge0$. Any matching matches at most $|N(S)|$ vertices of $S$, so leaves at least $d$ vertices of $X$ free for the best $S$:
$\nu\le|X|-d$. Add $d$ new vertices to $Y$, adjacent to every vertex of $X$. Now $|N'(S)|=|N(S)|+d\ge|S|$ for non-empty $S$, so Hall's theorem gives a matching covering $X$;
at most $d$ of its edges use new vertices, so the others form a matching of $G$ of size at least $|X|-d$.
\end{solution}

\begin{solution}{ex:edgecolour}
\idea A $\Delta$-regular bipartite multigraph splits into $\Delta$ perfect matchings, the colour classes.
\writeup Add isolated vertices to the smaller side so that $|X|=|Y|$. While some vertex has degree less than $\Delta$, there is one on each side: the degrees on each side add up to the number $m$ of edges, and if a vertex of $X$ has degree
less than $\Delta$, then $m<\Delta|X|=\Delta|Y|$, so some vertex of $Y$ has degree less than $\Delta$. Add an edge (possibly parallel) between two such
vertices. We reach a $\Delta$-regular bipartite multigraph containing $G$. By Corollary~\ref{cor:regular} its
edges split into $\Delta$ perfect matchings; colour each with its own colour and keep the edges of $G$. For the exams: students and examiners are the two
sides, exams the edges, and a colour is a time slot in which nobody has two exams. At least $\Delta$ slots are needed, and $\Delta$ suffice.
\end{solution}

\begin{solution}{ex:pathcover}
\idea A matching edge $u^-v^+$ says ``$v$ comes right after $u$ on its path''.
\writeup In a matching $M$ each vertex has at most one chosen successor and at most one chosen predecessor, so the chosen arcs form disjoint
directed paths and cycles; there are no cycles, as the graph is acyclic. Each matching edge joins two of the $n$ trivial paths, so the
vertices are covered by $n-|M|$ paths. Conversely, a cover by $q$ disjoint paths uses $n-q$ arcs, which form a matching; so $q\ge n-\nu$.
\end{solution}

\begin{solution}{ex:baseball}
\idea Integrality turns a flow into an outcome of the remaining games.
\writeup Let $k$ win all its remaining games, which is clearly best for it, and let $G=\sum_{\{i,j\}}g_{ij}$, the sum over pairs not containing $k$. An outcome of the
other games in which every team $i$ ends with at most $W$ wins gives a flow of value $G$: send from $\{i,j\}$ to $i$ the number of games $i$ wins against $j$;
team $i$ receives its number of new wins, at most $W-w_i$. Conversely, an integral flow of value $G$ (one exists if the maximum flow is $G$, by the
integrality theorem) tells how many of the $g_{ij}$ games each of $i$ and $j$ wins, and the capacities $W-w_i$ ensure that nobody passes $W$.
\end{solution}

\begin{solution}{ex:closure}
\idea A finite cut cannot cut an infinite arc, so it respects the requirements; its capacity counts lost profits and paid costs.
\writeup Let $Q=1+\sum_i|p_i|$ play the role of $+\infty$. For a cut $S$, let $C=S\setminus\{s\}$. If $C$ is allowed (no requirement leaves $C$), the arcs leaving $S$ are
$s\to i$ for $i\notin C$ with $p_i>0$, and $i\to t$ for $i\in C$ with $p_i<0$, so
\[
c(S)=\sum_{i\notin C,\,p_i>0}p_i+\sum_{i\in C,\,p_i<0}(-p_i)=\sum_{p_i>0}p_i-\sum_{i\in C}p_i.
\]
This is less than $Q$. If $C$ is not allowed, an arc of capacity $Q$ leaves $S$, so $c(S)\ge Q$, more than the cut $\{s\}$, of capacity at most $\sum_{p_i>0}p_i<Q$. So the
minimum cut is attained by an allowed $C$, and it minimises $\sum_{p_i>0}p_i-\mathrm{profit}(C)$.
\end{solution}

\begin{solution}{ex:birkhoffext}
\idea A $0$ entry of an average of non-negative matrices forces $0$ in each; an entry $1$ forces $1$ in each.
\writeup If $P=\frac12(B+C)$ with $P$ a permutation matrix and $B,C$ doubly stochastic, then where $P$ has $0$, $B$ and $C$ have non-negative entries adding up to $0$,
hence $0$; where $P$ has $1$, both entries are at most $1$ and average $1$, hence equal to $1$. So $B=C=P$. Conversely, if $A$ is not a permutation matrix,
Theorem~\ref{thm:birkhoff} writes $A=\lambda P+(1-\lambda)A'$ with $0<\lambda<1$ and $A'$ doubly stochastic, $A'\ne P$. For small $\varepsilon>0$, the matrices
$A\pm\varepsilon(P-A')$ lie on the segment from $A'$ to $P$, so they are doubly stochastic, they differ, and $A$ is their midpoint.
\end{solution}

\begin{solution}{ex:errera}
\idea Bridges are the only odd components that cost $S$ just one edge; with two of them the count still closes, by parity.
\writeup In Figure~\ref{fig:sylvester}, each blob is $K_4$ with one edge subdivided by a vertex joined to $s$; all degrees are $3$, and $S=\{s\}$ gives $o(G-S)=3>1$.
Now let $G$ be $3$-regular with at most two bridges, and $S$ a set of vertices. As in Corollary~\ref{cor:petersen}, $e(C,S)$ is odd for every odd component $C$ of
$G-S$, and if $e(C,S)=1$ that edge is a bridge, a different one for each $C$. So at most two odd components have $e(C,S)=1$, the others $e(C,S)\ge3$, and
\[
3\,o(G-S)-4\le\sum_Ce(C,S)\le3|S|.
\]
So $o(G-S)\le|S|+\frac43$, that is $o(G-S)\le|S|+1$. But $n$ is even ($3n$ is twice the number of edges), and $n\equiv|S|+o(G-S)\pmod 2$; so $o(G-S)\ne|S|+1$,
and $o(G-S)\le|S|$. Tutte's theorem concludes.
\end{solution}

\begin{solution}{ex:roommates}
\idea Everybody's last choice is $d$, and $a,b,c$ chase each other in a circle.
\writeup The roles of $a,b,c$ are cyclic ($a\to b\to c\to a$ in first choices), so it suffices to treat $d$ with $a$, the others following by rotation. Then $b$
and $c$ share. Student $c$ prefers $a$ to $b$, and $a$ prefers $c$ to $d$: $(a,c)$ blocks. So no pairing is stable. The mixed version is a marriage problem
with two proposers and two receivers, and Theorem~\ref{thm:galeshapley} gives a stable matching: bipartiteness is what saves stability,
as it saved Hall's condition.
\end{solution}

\begin{solution}{ex:stablelattice}
\idea A receiver claimed twice would have to be preferred both ways; a blocking pair for the new matching would block $M$ or $M'$.
\writeup Write $p\vee$ for the better of $M(p)$ and $M'(p)$ for $p$. Suppose $p\ne q$ both get $r$, with $r=M(p)$ and $r=M'(q)$ (it cannot come from the same matching
twice). Since $p$ chose $M(p)$, $p$ prefers $r$ to $M'(p)$ (strictly: $M'(p)\ne r=M'(q)$); stability of $M'$ forces $r$ to prefer $M'(r)=q$ to $p$. Likewise $q$ prefers $r$ to $M(q)$, and
stability of $M$ forces $r$ to prefer $M(r)=p$ to $q$. Contradiction: the new assignment is injective, hence a perfect matching.

Each receiver $r$ gets the worse of $M(r)$ and $M'(r)$. If they are equal there is nothing to prove. Otherwise let $p=M(r)\ne q=M'(r)$; the
receiver $r$ is now assigned to $p$ or to $q$, say to $p$ (the other case is symmetric). Then $p$ prefers $r=M(p)$ to $M'(p)$, and since $(p,r)$ does not block
$M'$, the receiver $r$ prefers $M'(r)=q$ to $p$: it got the worse of the two.

Stability: suppose that $(p,r)$ blocks the new matching. Then $p$ prefers $r$ to its new partner, hence to both $M(p)$ and $M'(p)$; and $r$ prefers $p$ to
its new partner, which is $M(r)$ or $M'(r)$. In the first case $(p,r)$ blocks $M$, in the second it blocks $M'$: a contradiction.
\end{solution}

\begin{solution}{ex:lying}
\idea A receiver's only power is to reject; rejecting its held proposer can set off a chain that brings back a better one.
\writeup Truthful run: $A$ proposes to $x$ (held), $B$ to $x$ ($x$ prefers $A$, rejects $B$), $C$ to $y$ (held), $B$ to $z$ (held). Outcome $Ax$, $Cy$, $Bz$: $x$ gets $A$, its second choice.
With $x$ reporting $C\,B\,A$: $A$ to $x$ (held); $B$ to $x$, who now rejects $A$; $C$ to $y$ (held); $A$ to $y$, who prefers $A$ and rejects $C$; $C$ to $x$, who rejects $B$; $B$ to $z$.
Outcome $Cx$, $Ay$, $Bz$: $x$ gets $C$, its true first choice. The rejection of $A$ sent $A$ to $y$, which released $C$.
\end{solution}

\section*{Sources and further reading}

\begin{itemize}
\item D.~Gale and L.~S.~Shapley, \emph{College admissions and the stability of marriage}, American Mathematical Monthly 69 (1962),
9--15; A.~E.~Roth, \emph{The evolution of the labor market for medical interns and residents: a case study in game theory},
Journal of Political Economy 92 (1984), 991--1016; A.~E.~Roth, \emph{Who Gets What and Why}, Houghton Mifflin Harcourt, 2015, a
lively account for everyone.
\item D.~G.~McVitie and L.~B.~Wilson, \emph{The stable marriage problem}, Communications of the ACM 14 (1971), 486--490; L.~E.~Dubins and
D.~A.~Freedman, \emph{Machiavelli and the Gale--Shapley algorithm}, American Mathematical Monthly 88 (1981), 485--494; D.~E.~Knuth,
\emph{Mariages stables}, Presses de l'Universit\'e de Montr\'eal, 1976 (English translation, AMS, 1997); B.~Pittel, \emph{On likely solutions
of a stable marriage problem}, Annals of Applied Probability 2 (1992), 358--401; D.~Gusfield and R.~W.~Irving, \emph{The Stable Marriage
Problem: Structure and Algorithms}, MIT Press, 1989.
\item C.~Berge, \emph{Two theorems in graph theory}, Proceedings of the National Academy of Sciences 43 (1957), 842--844; P.~Hall,
\emph{On representatives of subsets}, Journal of the London Mathematical Society 10 (1935), 26--30; P.~R.~Halmos and H.~E.~Vaughan,
\emph{The marriage problem}, American Journal of Mathematics 72 (1950), 214--215.
\item L.~R.~Ford and D.~R.~Fulkerson, \emph{Maximal flow through a network}, Canadian Journal of Mathematics 8 (1956), 399--404;
A.~Schrijver, \emph{On the history of the transportation and maximum flow problems}, Mathematical Programming 91 (2002), 437--445,
which tells the story of the Harris--Ross report.
\item G.~Birkhoff, \emph{Tres observaciones sobre el \'algebra lineal}, Universidad Nacional de Tucum\'an, Revista Serie A 5 (1946), 147--151.
\item W.~T.~Tutte, \emph{The factorization of linear graphs}, Journal of the London Mathematical Society 22 (1947), 107--111; L.~Lov\'asz,
\emph{Three short proofs in graph theory}, Journal of Combinatorial Theory B 19 (1975), 269--271; J.~Petersen, \emph{Die Theorie der
regul\"aren Graphs}, Acta Mathematica 15 (1891), 193--220; J.~Edmonds, \emph{Paths, trees, and flowers}, Canadian Journal of Mathematics
17 (1965), 449--467.
\item L.~Lov\'asz and M.~D.~Plummer, \emph{Matching Theory}, North-Holland, 1986 (reprinted by the AMS, 2009); A.~Schrijver,
\emph{Combinatorial Optimization: Polyhedra and Efficiency}, Springer, 2003, the encyclopaedia of the subject.
\end{itemize}
